\scanpage{60}
\mspage{59}

\uline{\unsure{Def|Déf}} Un atelier est dit ``ensembliste \unsure{spacieux|spécieux}'' s'il satisfait la condition (i). Idem pour ``ensembliste polyédral'', en remplaçant $\mathrm{Fig}(L)$ par $\mathrm{Fig}\,\mathrm{pol}(L)$. On a les \deleted{conditions} notions correspondantes pour les préateliers.

% margin note written sideways; slanted upward
\margin{NB. Il faut voir les figures dans un tel atelier comme étant des figures ensemblistes dans $L$, munies de \unsure{structures suppl\textsuperscript{taires}|structures supplémentaires} (\unsure{tout} comme les \unsure{2-figures|S-figures} en $L$, exemple plus haut).}

Un atelier est dit ``quasi-ensembliste \unsure{spacieux|spécieux}'' s'il satisfait (ii). Notions correspondantes pour le cas ``quasi-ensembliste polyédral'', ou le cas des préateliers.

Les préateliers quasi-ensemblistes (indifféremment \unsure{spacieux|spécieux} ou polyédraux) correspondent aux $(M,\leq,\ll)$ satisfaisant les conditions suivantes (à déduire de $\unsure{\mathrm{Prat}|\mathrm{Prot}}\,1'-\mathrm{Prat}\,6'$) :

\uline{Prat quas} 1
\[
X\leq Y\ \Rightarrow\ X\ll Y
\]

% slanted margin labels, with arrows toward the numbered conditions
\margin{densité\newline divisibilité $\Longrightarrow$}
\uline{Prat quas} 2
\[
\forall X\in M,\quad\operatorname{omb}(X)^\circ\ne\varnothing\quad\text{i.e.}\quad\exists x\in L,\quad x\overset{\circ}{\ll}X
\]

\margin{\unsure{Aut. minimale} $\longrightarrow$}
\uline{Prat quas} 3
\[
\forall X\in M,\ x\in L,\ x\ll X,\quad\exists X'\in M,\quad x\overset{\circ}{\ll}X'\leq X
\]

\margin{raffinement induit $\longrightarrow$}
\uline{Prat quas} 4 Soit
\diagram{\begin{tikzcd}[ampersand replacement=\&,column sep=small,row sep=small]
Y \arrow[r,phantom,"\ll" description] \& X \arrow[d,phantom,"\unsure{\vee|\geq}" description] \\
\& X'
\end{tikzcd}}
\deleted{Alors} et $Y\deleted{\unsure{X'}}_{X'}=\widetilde Y\cap X'$.

Alors
\[
\operatorname{omb}(Y)\cap\operatorname{omb}(X')=\operatorname{omb}(\unsure{Y'|Y_{X'}})\left(\overset{\mathrm{def}}{=}\bigcup_{Z\in Y_{X'}}\operatorname{omb}(Z)\right),
\]
i.e. si
\[
x\in L,\quad x\ll Y,X',\qquad\exists Y'\text{ avec }Y'\leq Y,\ Y'\ll X',\ x\ll Y'.
\]
De plus, si $Y\leq X$, OPS $Y'\leq X'$ (i.e. on peut remplacer $Y_{X'}$ par $\widetilde Y\cap\widetilde X'$ : -- les deux \unsure{seraient|seront} d'ailleurs égaux).

Il n'y a plus lieu en effet d'introduire $\parallel_L$, vu qu'on la connaît : \deleted{\unsure{c.à.d.}} $(x\parallel y)\Longleftrightarrow(x\ne y)$.
